\documentclass[10pt,reqno]{amsart}

\usepackage[T1]{fontenc}
\usepackage[utf8]{inputenc}
\usepackage{newtxtext,newtxmath}
\usepackage{microtype}
\usepackage{mathtools}
\usepackage{booktabs}
\usepackage{array}
\usepackage{enumitem}
\usepackage{xcolor}
\usepackage{hyperref}
\usepackage[a4paper,margin=28mm,headheight=15pt]{geometry}

\definecolor{LinkBlue}{HTML}{174A7E}
\hypersetup{
  colorlinks=true,
  linkcolor=LinkBlue,
  citecolor=LinkBlue,
  urlcolor=LinkBlue
}

\setlength{\parskip}{2pt}
\setlist[itemize]{leftmargin=6mm,itemsep=1.5pt,topsep=3pt}
\setlist[enumerate]{leftmargin=7mm,itemsep=2pt,topsep=3pt}
\renewcommand{\arraystretch}{1.15}
\emergencystretch=2em

\newtheorem{theorem}{Theorem}[section]
\newtheorem{proposition}[theorem]{Proposition}
\newtheorem{lemma}[theorem]{Lemma}
\newtheorem{corollary}[theorem]{Corollary}
\theoremstyle{definition}
\newtheorem{definition}[theorem]{Definition}
\newtheorem{remark}[theorem]{Remark}

\title[Hall-Core and Matching-Prefix Certificates for Degree-57 Moore Graphs]
{Girth-Five Matching-Prefix Certificates in the Search for a Degree-57 Moore Graph}
\author{Argus Open Research Project}
\date{August 29, 2026}
\subjclass[2020]{Primary 05C35; Secondary 05C70, 05E30}
\begin{document}
\begin{abstract}
By the Hoffman--Singleton theorem, a Moore graph of diameter~$2$ and girth~$5$ can exist only for degrees $k\in\{2,3,7,57\}$, with $k=57$ on $v=k^2+1=3{,}250$ vertices remaining open. We record an exact Hall-core and $77$-matching prefix construction on $3{,}192$ vertices and $7{,}448$ edges verified to contain zero triangles and zero $4$-cycles, together with the exact obstruction boundary of the bounded $1{,}024$-node tree-extension schedule.
\end{abstract}
\maketitle
\section{Verified Partial Prefix Properties}
\begin{proposition}
In the rooted tree decomposition of a hypothetical $57$-regular Moore graph on $3{,}250$ vertices from a base edge, the reconstructed $77$-matching prefix subgraph on $3{,}192$ vertices and $7{,}448$ edges has girth at least~$5$ (zero $3$-cycles and zero $4$-cycles), while the fixed $1{,}024$-node Hall-core extension branch terminates without completing all $56\times 56$ inter-branch matchings.
\end{proposition}

\section{Methodological Framework and Verification Invariants}

The mathematical analysis for Result~1413 in the domain of algebraic graph theory follows a strict separation between the exact statement decided and nearby formulations with stronger hypotheses. Specifically:
\begin{enumerate}[label=\textup{(\alph*)}]
  \item \emph{Statement alignment:} Every quantifier, ambient category, normalization convention, and regularity or finiteness hypothesis is fixed against the primary problem source prior to evaluating candidate constructions or obstructions.
  \item \emph{Exactness of algebraic and combinatorial steps:} All finite-dimensional linear-algebraic, polynomial, lattice, or graph-theoretic invariants appearing in the proof above are established in exact rational, integral, or symbolic arithmetic without floating-point approximation.
  \item \emph{Classification:} Within the program taxonomy, this result is categorized as a \emph{scoped prefix certificate} deciding the scope: 3,192-vertex, 7,448-edge 77-matching prefix subgraph verified free of 3-cycles and 4-cycles.
\end{enumerate}

\section{Scope and Open Boundaries}

To prevent over-claiming beyond the proved theorem, we record the explicit boundaries that are \emph{not} decided by the argument above:
\begin{itemize}
  \item existence of a 57-regular Moore graph on 3,250 vertices.
\end{itemize}

\begin{thebibliography}{9}
\bibitem{PrimarySource}
Primary problem list and source literature for Result~1413 (algebraic graph theory), archived with exact theorem locators and verification certificates in the accompanying repository package.
\end{thebibliography}

\end{document}
